\chapter{Matrix Representation of Linear Maps}
\label{ch:i07-matrix-representation-of-linear-maps}

Matrices are coordinate representations of linear maps. The matrix is not the transformation itself: it depends on chosen bases in the domain and codomain. Once this distinction is kept visible, matrix multiplication, change of basis, rank and similarity all acquire transparent conceptual meanings.

\section*{Learning goals}
After this chapter, the reader should be able to:
\begin{itemize}
\item construct the matrix of a linear map from its action on a basis;
\item recover the action of a linear map from its matrix and coordinate conventions;
\item compute the matrix of a composition and explain why matrix multiplication appears in that order;
\item change the matrix of an operator when the basis changes;
\item identify the column space and null space of a matrix as image and kernel data of the associated map;
\item check a matrix representation by applying it to basis vectors and comparing the resulting coordinates;
\end{itemize}
\section*{Conceptual roadmap}
The chapter is organized around a repeated pattern:
\[
\boxed{\text{definition}\;\longrightarrow\;\text{structural theorem}\;\longrightarrow\;\text{coordinate calculation}\;\longrightarrow\;\text{application}.}
\]
Whenever possible, we separate the invariant mathematical object from a chosen coordinate representation.

\section{Matrix of a linear map}
For bases $B$ of $V$ and $C$ of $W$, the matrix $[T]_{C\leftarrow B}$ is defined by $[T(v)]_C=[T]_{C\leftarrow B}[v]_B$.

\section{Column rule}
The $j$th column of $[T]_{C\leftarrow B}$ is $[T(v_j)]_C$.

\section{Composition and multiplication}
Coordinates turn composition into matrix multiplication: $[S\circ T]=[S][T]$ with compatible bases.

\section{Identity and inverses}
The identity has the identity matrix in the same basis. An isomorphism is represented by an invertible matrix.

\section{Change of basis for operators}
For an endomorphism, changing from basis $B$ to $C$ gives $[T]_C=P_{C\leftarrow B}[T]_B P_{B\leftarrow C}$. Thus representations of the same operator are similar.

\section{Rank and row reduction}
The rank of a matrix equals the dimension of the image of its associated linear map. Row operations preserve rank.

\section{Core structural results}

\begin{theorem}[Matrix composition rule]\label{thm:i07-01}
With compatible bases, $[S\circ T]=[S][T]$.
\begin{proof}
Apply both sides to an arbitrary coordinate vector and use the defining representation identity twice.
\end{proof}
\end{theorem}

\begin{theorem}[Similarity under basis change]\label{thm:i07-02}
If $A=[T]_B$ and $P=P_{C\leftarrow B}$, then $[T]_C=PAP^{-1}$.
\begin{proof}
Use $[v]_C=P[v]_B$ and express the coordinate identity for $T(v)$ in the two bases.
\end{proof}
\end{theorem}

\section{Worked examples}

\begin{example}[Derivative matrix]\label{ex:i07-worked-01}
On $\mathcal P_2$ in the monomial basis, differentiation is represented by $\begin{pmatrix}0&1&0\\0&0&2\\0&0&0\end{pmatrix}$.
\end{example}

\begin{example}[Rotation]\label{ex:i07-worked-02}
A planar rotation through angle $\theta$ has standard matrix $\begin{pmatrix}\cos\theta&-\sin\theta\\\sin\theta&\cos\theta\end{pmatrix}$.
\end{example}

\begin{example}[Projection]\label{ex:i07-worked-03}
Projection onto the first coordinate axis has matrix $\begin{pmatrix}1&0\\0&0\end{pmatrix}$.
\end{example}

\section{Diagnostic checklist}
Before moving on, it is useful to distinguish four levels of understanding:
\begin{enumerate}
\item recognize the definitions in unfamiliar examples;
\item prove the core structural statements without coordinates when possible;
\item translate the same statement into a matrix or coordinate computation;
\item know which hypotheses are essential and produce a counterexample when one is removed.
\end{enumerate}

% BEGIN VOLUME01 SOLVED DOSSIERS
\section{Solved dossiers}
These dossiers are longer worked problems. They complement the shorter exercise--hint--solution triads by combining several ideas in one argument.

\begin{problem}[Matrix from basis images]\label{prob:i07-dossier-01}
Find the standard matrix of $T(x,y,z)=(x+z,2y,x-z)$.
\end{problem}
\begin{solution}
The columns are $T(e_1)=(1,0,1)$, $T(e_2)=(0,2,0)$, and $T(e_3)=(1,0,-1)$, giving $\begin{pmatrix}1&0&1\\0&2&0\\1&0&-1\end{pmatrix}$.
\end{solution}

\begin{problem}[Composition matrix]\label{prob:i07-dossier-02}
Let $A=\begin{pmatrix}1&1\\0&1\end{pmatrix}$ and $B=\begin{pmatrix}0&-1\\1&0\end{pmatrix}$. Which matrix represents first applying $A$ and then $B$?
\end{problem}
\begin{solution}
Column vectors are acted on from the left, so first $A$ then $B$ is represented by $BA$. Matrix multiplication encodes composition in reverse verbal order.
\end{solution}

\begin{problem}[Similarity from basis change]\label{prob:i07-dossier-03}
If $[T]_B=A$ and $P=P_{C\leftarrow B}$, derive $[T]_C=PAP^{-1}$.
\end{problem}
\begin{solution}
For $x_C=Px_B$, we have $[T(v)]_C=P A x_B=PAP^{-1}x_C$. Since this holds for every coordinate column, the $C$-matrix is $PAP^{-1}$.
\end{solution}

\begin{problem}[Derivative matrix]\label{prob:i07-dossier-04}
Write the matrix of $D:\mathcal P_3\to\mathcal P_2$ in monomial bases.
\end{problem}
\begin{solution}
$D(1)=0$, $D(x)=1$, $D(x^2)=2x$, $D(x^3)=3x^2$. Thus the columns are $(0,0,0)^T,(1,0,0)^T,(0,2,0)^T,(0,0,3)^T$.
\end{solution}

\begin{problem}[Rank from columns]\label{prob:i07-dossier-05}
For $A=\begin{pmatrix}1&2&3\\2&4&6\end{pmatrix}$ find the rank.
\end{problem}
\begin{solution}
Every column is a scalar multiple of $(1,2)^T$, and the matrix is nonzero. Hence the column space is one-dimensional and rank is $1$.
\end{solution}

\begin{problem}[Elementary matrix meaning]\label{prob:i07-dossier-06}
What linear map does left multiplication by the elementary matrix that swaps rows $1$ and $2$ represent?
\end{problem}
\begin{solution}
It composes a matrix-valued coordinate representation with the coordinate permutation swapping the first two output coordinates. Algebraically it performs the corresponding row operation.
\end{solution}

\begin{problem}[Inverse matrix as inverse map]\label{prob:i07-dossier-07}
Why does $A^{-1}$ represent $T^{-1}$ when $A$ represents an isomorphism $T$ in fixed bases?
\end{problem}
\begin{solution}
The composition $T^{-1}T=I$ gives $[T^{-1}]A=I$ by the composition rule, so $[T^{-1}]=A^{-1}$.
\end{solution}

\begin{problem}[Rectangular matrix]\label{prob:i07-dossier-08}
What does a $3\times2$ matrix represent geometrically?
\end{problem}
\begin{solution}
It represents a linear map from a two-dimensional coordinate space to a three-dimensional coordinate space. Its two columns are the images of the two domain basis vectors expressed in the codomain basis.
\end{solution}

\begin{problem}[Same operator, different matrices]\label{prob:i07-dossier-09}
Can two different matrices represent the same linear operator?
\end{problem}
\begin{solution}
Yes, if different bases are used. The matrices are related by similarity for an endomorphism. With the basis fixed, however, the representing matrix is unique.
\end{solution}

\begin{problem}[Projection matrix]\label{prob:i07-dossier-10}
Find the standard matrix for projection onto $\operatorname{span}(e_1,e_3)$ in $\mathbb R^3$.
\end{problem}
\begin{solution}
The map sends $(x,y,z)$ to $(x,0,z)$, so the matrix is $\operatorname{diag}(1,0,1)$.
\end{solution}

\begin{problem}[Column space and image]\label{prob:i07-dossier-11}
Explain why the column space of a matrix is the coordinate image of its linear map.
\end{problem}
\begin{solution}
Every output is $Ax=\sum_jx_j a_j$, a linear combination of the columns $a_j$. Conversely every such combination occurs for the coefficient vector $x$, so the image equals the column span.
\end{solution}

\begin{problem}[Matrix synthesis]\label{prob:i07-dossier-12}
Why is similarity an equivalence relation on square matrices?
\end{problem}
\begin{solution}
Reflexivity uses $I$, symmetry uses $P^{-1}$, and transitivity follows because $B=PAP^{-1}$ and $C=QBQ^{-1}$ imply $C=(QP)A(QP)^{-1}$. It therefore captures 'same operator in different bases'.
\end{solution}

% END VOLUME01 SOLVED DOSSIERS

\section{Exercises with complete solutions}

\begin{exercise}\label{exr:i07-01}
Find the matrix of $T(x,y)=(x+2y,3x-y)$ in the standard basis.
\end{exercise}
\begin{hint}
Identify the definition or structural theorem in this chapter that directly controls the question, then reduce the calculation to that statement.
\end{hint}
\begin{solution}
The columns are $T(e_1)=(1,3)$ and $T(e_2)=(2,-1)$, so $\begin{pmatrix}1&2\\3&-1\end{pmatrix}$.
\end{solution}

\begin{exercise}\label{exr:i07-02}
Why do columns record images of basis vectors?
\end{exercise}
\begin{hint}
Identify the definition or structural theorem in this chapter that directly controls the question, then reduce the calculation to that statement.
\end{hint}
\begin{solution}
The coordinate vector of the $j$th basis vector is $e_j$, and multiplying by a matrix selects its $j$th column.
\end{solution}

\begin{exercise}\label{exr:i07-03}
Prove that the matrix of an inverse map is the inverse matrix.
\end{exercise}
\begin{hint}
Identify the definition or structural theorem in this chapter that directly controls the question, then reduce the calculation to that statement.
\end{hint}
\begin{solution}
From $T^{-1}T=I$, the composition rule gives $[T^{-1}][T]=I$.
\end{solution}

\begin{exercise}\label{exr:i07-04}
Compute the rank of $\begin{pmatrix}1&2\\2&4\end{pmatrix}$.
\end{exercise}
\begin{hint}
Identify the definition or structural theorem in this chapter that directly controls the question, then reduce the calculation to that statement.
\end{hint}
\begin{solution}
The second row and second column are scalar multiples of the first, while the matrix is nonzero, so rank is $1$.
\end{solution}

\begin{exercise}\label{exr:i07-05}
Show that similar matrices have the same rank.
\end{exercise}
\begin{hint}
Identify the definition or structural theorem in this chapter that directly controls the question, then reduce the calculation to that statement.
\end{hint}
\begin{solution}
Left and right multiplication by invertible matrices correspond to isomorphisms and do not change image dimension.
\end{solution}

\begin{exercise}\label{exr:i07-06}
Write the standard matrix of reflection across the $x$-axis.
\end{exercise}
\begin{hint}
Identify the definition or structural theorem in this chapter that directly controls the question, then reduce the calculation to that statement.
\end{hint}
\begin{solution}
$\begin{pmatrix}1&0\\0&-1\end{pmatrix}$.
\end{solution}

\begin{exercise}\label{exr:i07-07}
Explain why matrix multiplication is generally noncommutative.
\end{exercise}
\begin{hint}
Identify the definition or structural theorem in this chapter that directly controls the question, then reduce the calculation to that statement.
\end{hint}
\begin{solution}
Composition of transformations is order-sensitive; for example, a projection followed by a rotation differs from the reverse order.
\end{solution}

\begin{exercise}\label{exr:i07-08}
Relate elementary row operations to invertible matrices.
\end{exercise}
\begin{hint}
Identify the definition or structural theorem in this chapter that directly controls the question, then reduce the calculation to that statement.
\end{hint}
\begin{solution}
Each elementary row operation is left multiplication by an elementary invertible matrix.
\end{solution}

\section*{Chapter summary}
Matrices are coordinate representations of linear maps. The matrix is not the transformation itself: it depends on chosen bases in the domain and codomain. Once this distinction is kept visible, matrix multiplication, change of basis, rank and similarity all acquire transparent conceptual meanings.
The important habit is to move fluently between invariant statements and concrete coordinates while remembering which conclusions depend on finite dimensionality, the scalar field, or the inner product.
